\FifthArticle{2}
The entire domain of complex quantities, in which the values of $x$ satisfying the equation $X=0$ are to be sought, is an infinite domain of two dimensions. For if such a value is put in the form $x=t+iu$, where $i$ always denotes the imaginary unit $\surd{-1}$, all real values from $-\infty$ to $+\infty$ are admissible for $t$ and $u$. We must first delimit within this infinite domain a finite one outside which no root of the given equation $X=0$ can possibly lie. This can be done in more than one way; the following seems best suited to our purpose.

Instead of the form $t+iu$, use
\[
x=r(\cos\rho+i\sin\rho)
\]
To cover the entire infinite domain of complex quantities, $r$ must then range over all positive values from $0$ to $+\infty$, and $\rho$ from $0$ to $360^0$, or, what
\SourceLeaf{083}
is the same, from any initial value to a final value greater by $360^0$.

To obtain a bound for $r$ beyond which no value can correspond to a root of the equation $X=0$, I first put the coefficients of the individual terms of $X$ into a form similar to that of $x$, namely
\[
\begin{array}{l}
A=a(\cos\alpha+i\sin\alpha)\\
B=b(\cos\beta+i\sin\beta)\\
C=c(\cos\gamma+i\sin\gamma)\quad\mbox{etc.}
\end{array}
\]
Here $a,b,c$ are to denote fixed positive quantities, except that one or another of them may also be $=0$. I then consider the equation
\[
r^n-\surd2\,.(ar^{n-1}+br^{n-2}+cr^{n-3}+\mbox{etc.})=0
\]
which, as is readily seen, has a positive root, and indeed, by \textsc{Harriot}'s theorem, only one such root. Let $R$ be this root. It is then evident that for every positive value of $r$ greater than $R$, the value of
$r^n-\surd2\,.(ar^{n-1}+br^{n-2}+cr^{n-3}+\mbox{etc.})$
is positive, and the same also holds for the function
\[
nr^n-\surd2\,.((n-1)ar^{n-1}+(n-2)br^{n-2}+(n-3)cr^{n-3}+\mbox{etc.})
\]
since it exceeds $n$ times the first function by
\[
\surd2\,.(ar^{n-1}+2br^{n-2}+3cr^{n-3}+\mbox{etc.})
\]
and thus by a positive difference.

\FifthArticle{3}
I now assert that the quantity $R$ can serve as a bound for the values of $r$ of the kind required in the preceding article. The proposition is to be proved as follows.

\Needspace{10\baselineskip}In general I put $X=T+iU$, where $T$ and $U$ of course denote real quantities. In fact,
\[
\begin{array}{r@{}l}
T={}&r^n\cos n\rho+ar^{n-1}\cos((n-1)\rho+\alpha)+br^{n-2}\cos((n-2)\rho+\beta)\\
&{}+cr^{n-3}\cos((n-3)\rho+\gamma)+\mbox{etc.}
\end{array}
\]
\SourceLeaf{084}
\[
\begin{array}{r@{}l}
U={}&r^n\sin n\rho+ar^{n-1}\sin((n-1)\rho+\alpha)+br^{n-2}\sin((n-2)\rho+\beta)\\
&{}+cr^{n-3}\sin((n-3)\rho+\gamma)+\mbox{etc.}
\end{array}
\]
\Needspace{15\baselineskip}It is easy to see that when any positive value greater than $R$ is chosen for $r$, $T$ necessarily has the same sign as $\cos n\rho$ whenever the absolute value of this cosine is not less than $\surd\frac12$. For this it is enough to put $T$ into the form
\[
\begin{array}{r@{}l}
\pm T={}&\surd\frac12\,.r^n-ar^{n-1}-br^{n-2}-cr^{n-3}-\mbox{etc.}\\
&{}+[\pm\cos n\rho-\surd\frac12]r^n\\
&{}+[1\pm\cos((n-1)\rho+\alpha)]ar^{n-1}\\
&{}+[1\pm\cos((n-2)\rho+\beta)]br^{n-2}\\
&{}+[1\pm\cos((n-3)\rho+\gamma)]cr^{n-3}\\
&{}+\mbox{etc.}
\end{array}
\]
The upper signs are to apply when $\cos n\rho$ is positive, and the lower signs when it is negative. The first part of the expression on the right is positive by the proposition of the preceding article, while all the remaining parts are at least nonnegative. In much the same way, by writing $U$ in place of $T$ in the above formula and replacing cosines by sines throughout, it follows that under the same assumption about $r$, $U$ always has the same sign as $\sin n\rho$ whenever the absolute value of this sine is not less than $\surd\frac12$. Thus, in every case, at least one of the two quantities $T,U$ has a predetermined positive or negative sign. Consequently, for no value of $\rho$ can the function satisfy $X=0$. Q.~E.~D.

\Needspace{16\baselineskip}\FifthArticle{4}
To clarify further the behaviour of $T$ and $U$ with regard to their signs and changes of sign, for a fixed value of $r$ exceeding $R$, let $\rho$ run through all values between two limits differing by $360^0$. Instead of $0$ and $360^0$, however, taking for brevity
\[
\frac{45^0}{n}=\omega
\]
\SourceLeaf{085}
the limits $-\omega$ and $(8n-1)\omega$ are to be chosen. I divide the entire range into $4n$ equal parts, the first extending from $-\omega$ to $\omega$, the second from $\omega$ to $3\omega$, the third from $3\omega$ to $5\omega$, and so on. First the values of the derivatives $\frac{\mathrm{d}T}{\mathrm{d}\rho}$, $\frac{\mathrm{d}U}{\mathrm{d}\rho}$ must also be considered. They are given by
\[
\begin{array}{rl}
\frac{\mathrm{d}T}{\mathrm{d}\rho}
&=-nr^n\sin n\rho-(n-1)ar^{n-1}\sin\bigl((n-1)\rho+\alpha\bigr)\\
&\qquad -(n-2)br^{n-2}\sin\bigl((n-2)\rho+\beta)\bigr)\\
&\qquad -(n-3)cr^{n-3}\sin\bigl((n-3)\rho+\gamma\bigr)-\mathrm{etc.}\\
\frac{\mathrm{d}U}{\mathrm{d}\rho}
&=nr^n\cos n\rho+(n-1)ar^{n-1}\cos\bigl((n-1)\rho+\alpha\bigr)\\
&\qquad +(n-2)br^{n-2}\cos\bigl((n-2)\rho+\beta)\bigr)\\
&\qquad +(n-3)cr^{n-3}\cos\bigl((n-3)\rho+\gamma\bigr)+\mathrm{etc.}
\end{array}
\]
\Needspace{6\baselineskip}By reasoning similar to that of the preceding article, and using the proposition at the end of Article 2, it is easy to see from these formulas that $\frac{\mathrm{d}T}{\mathrm{d}\rho}$ always has the opposite sign to $\sin n\rho$ whenever the absolute value of this sine is not less than $\surd\frac{1}{2}$, whereas $\frac{\mathrm{d}U}{\mathrm{d}\rho}$ always has the same sign as $\cos n\rho$ whenever the absolute value of this cosine is not less than $\surd\frac{1}{2}$. The following conclusions result.

\Needspace{14\baselineskip}In the first interval, that is, from $\rho=-\omega$ to $\rho=+\omega$, $T$ is always positive, whereas $U$ is negative at the initial value and positive at the final value. It is therefore certainly $=0$ once in between, and \emph{only} once, since $\frac{\mathrm{d}U}{\mathrm{d}\rho}$ is positive throughout the interval.

In the second interval, $U$ is always positive, while $T$ is positive at the beginning and negative at the end. In between, $T=0$ holds once, and only once, since $\frac{\mathrm{d}T}{\mathrm{d}\rho}$ is negative throughout the interval.

In the third interval, $T$ is always negative, while $U$ undergoes a change of sign, so that $U=0$ holds \emph{once}.

In the fourth interval, $U$ is always negative, and $T$ is $=0$ \emph{once}.

In the following intervals these relations repeat in the same order, so that the fifth corresponds to the first, the sixth to the second, and so on.

